\begin{figure}[H]%[ht]
	\centering
	\begin{equation*}
		\begin{array}{c|cccccc|c}
			& x_1 & x_2 & \dots & x_j & \dots & x_n & \\
			\hline
			y_1 & a_{11} & a_{12} & \dots & a_{1j} & \dots & a_{1n} & b_1 \\
			y_2 & a_{21} & a_{22} & \dots & a_{2j} & \dots & a_{2n} & b_2 \\
			\vdots & \vdots & \vdots & \ddots & \vdots & \ddots & \vdots & \vdots \\
			y_i & a_{i1} & a_{i2} & \dots & a_{ij} & \dots & a_{in} & b_i \\
			\vdots & \vdots & \vdots & \ddots & \vdots & \ddots & \vdots & \vdots \\
			y_m & a_{m1} & a_{m2} & \dots & a_{mj} & \dots & a_{mn} & b_m \\
			\hline
			& c_1 & c_2 & \dots & c_j & \dots & c_n &
		\end{array}
	\end{equation*}
	\caption{Rappresentazione schematica della dualit\`a nei sistemi lineari.}
	\label{fig:duale}
\end{figure}
